% ==============================================================================
% SECCION VI: CONCLUSIONES
% ==============================================================================
\section{Conclusiones}\label{sec:conclusiones}

\subsection{Respuesta a la Pregunta de Investigación y Síntesis de Hallazgos}\label{subsec:resp_conclusiones}

En respuesta directa y explícita a la pregunta de investigación planteada en la Sección~\ref{sec:introduccion}, los resultados experimentales y analíticos de este estudio demuestran que \textbf{sí es factible} lograr levantamiento autónomo rápido y estabilización invertida continua y robusta a \qty{1.0}{\kilo\hertz} sobre una plataforma impresa en PLA con puente~H L298N y encoders magnéticos AS5600, venciendo sus no linealidades severas \textbf{sin recurrir a sensores ópticos ni puentes MOSFET}. El éxito sobre este hardware económico no se alcanza ajustando reguladores lineales sobre modelos ideales, sino mediante la co-ingeniería acoplada de cuatro pilares: \textit{(i)}~manufactura modular libre de colisiones con desacoplamiento inercial $80:1$ y sensado acimutal previo a la caja reductora; \textit{(ii)}~instrumentación isócrona determinista con doble bus Soft-$\text{I}^2\text{C}$ dedicado y modulación PWM ultrasónica; \textit{(iii)}~identificación empírica en banco de la zona muerta asimétrica, la planta motor--reductor y la fricción mixta; y \textit{(iv)}~la arquitectura híbrida \texttt{M4 Ultra-Silk}, que integra levantamiento resonante en cuadratura acotado en carrera, lazo dual paralelo sin acción integral en el péndulo y con signos de fase no mínima en el brazo, compensación continua por tramos del actuador, cascada de filtrado pasa-bajos y \textit{Dual-Notch} adaptativo de 64~bits, y retorno servo con reinicio de Clegg y freno activo bajo supervisión en tiempo real de un gemelo digital~3D.

En concordancia con esta respuesta, la investigación verifica el cumplimiento cuantitativo de los cuatro objetivos específicos propuestos:

\begin{enumerate}
    \item \textbf{Plataforma modular DfAM e instrumentación determinista a \qty{1.0}{\kilo\hertz}:} La estructura del péndulo impresa en PLA sin material de soporte ($\text{SUP}=0$) en cuatro placas consumió \qty{178.13}{\gram} (\qty{59.72}{\meter} de filamento) en 15h~10m de fabricación y superó la validación booleana sin colisiones ($0/23$ interferencias en el ensamble completo y $0/6$ en el subconjunto motor--encoder). El firmware desplegado sobre STM32F446RE a \qty{180}{\mega\hertz} ejecutó el lazo isócrono \texttt{TIM6} a \qty{1.0}{\kilo\hertz}, la portadora PWM en \texttt{TIM1} a \qty{20.0}{\kilo\hertz} y la adquisición simultánea por doble bus Soft-$\text{I}^2\text{C}$ a una frecuencia cercana a \qty{450}{\kilo\hertz} con tiempo de lectura inferior a \qty{45}{\micro\second} por sensor, multiplicando ochenta veces la resolución efectiva del brazo al ubicar el AS5600 base en el eje del rotor antes de la reductora $80:1$.

    \item \textbf{Caracterización e identificación experimental en banco real:} A partir de una muestra total de $N=3\,726$ registros adquiridos a \qty{100}{\hertz} por el enlace serial \texttt{COM13} a \qty{115200}{\baud}, se identificó la zona muerta asimétrica del conjunto L298N--reductor en $+10.0\%$ (antihorario) y $-21.0\%$ (horario) de ciclo PWM, compensada de forma continua por tramos~\eqref{eq:comp_l298n} en $+9.2\%$ y $-16.0\%$ con tensión de rodilla $V_k=\qty{0.65}{\volt}$. Además, se caracterizó la función de transferencia del actuador $G_v(s)$~\eqref{eq:gv_identificada} con $K_v=\qty{3.3720}{(\radian\per\second)\per\volt}$ y $\tau_m=\qty{0.1175}{\second}$ ($R^2=98.51\%$, $\mathrm{NRMSE}=87.79\%$), junto con la dinámica pendular acoplada ($\omega_n=\qty{8.8484}{\radian\per\second}$, $f_n=\qty{1.4083}{\hertz}$, $\zeta=0.06488$, $K_{\text{acople}}=0.8737$) y el modelo mixto viscoso--Coulomb~\eqref{eq:friccion_mixta} con coeficiente de determinación $R^2=99.64\%$.

    \item \textbf{Arquitectura híbrida \texttt{M4 Ultra-Silk} y gemelo digital~3D:} La síntesis física en lazo dual paralelo~\eqref{eq:m4_dual} suprimió la divergencia causada por la competencia integral SIMO al imponer ganancia integral nula en el lazo pendular y signos invertidos de fase no mínima en el brazo, certificando estabilidad asintótica modal con radio espectral discreto $\rho=0.99934$~\eqref{eq:polos_m4}, estrictamente menor que la unidad. La cascada doble pasa-bajos IIR~\eqref{eq:lpf_num_z} combinada con el bloque \textit{Dual-Notch} en Forma Directa~II Transpuesta (TDF-II) de 64~bits ---muesca fija en \qty{11.0}{\hertz}~\eqref{eq:hz_n11} y rama adaptativa en \qty{6.2}{\hertz}~\eqref{eq:hz_n6} modulada según la severidad de perturbación~\eqref{eq:mezcla_dist}--- atenuó el ruido derivativo y el ciclo límite de holgura con un retardo de fase de apenas \qty{-12.87}{\degree} en $f_n$. El lazo servo \texttt{CHAR5B} con reinicio de Clegg~\eqref{eq:clegg_pos} y freno activo logró un sobreimpulso $M_p=0.29\%$ en simulación y retorno físico \texttt{STOPHOME} en \qty{-0.25}{\degree}, mientras el canal \texttt{DMA1} circular sostuvo la telemetría CSV a \qty{100}{\hertz} y alimentó en tiempo real la animación cinemática a \qty{30}{\hertz} (periodo de \qty{0.033}{\second}) del gemelo digital~3D en MATLAB (\texttt{app\_telemetria\_sistema\_\allowbreak furuta\_matlab.m}).

    \item \textbf{Validación cuantitativa cruzada entre modelo, simulación y hardware:} Los parámetros inerciales y electromecánicos identificados concordaron dentro de un error relativo inferior al \qty{3}{\percent} respecto al diseño CAD, explicándose el incremento de amortiguamiento y la asimetría de arranque por los sellos del rodamiento distal y la topología Darlington. En la validación física continua de \qty{15.001}{\second} a \qty{12.0}{\volt}, el prototipo completó el levantamiento autónomo en \qty{0.980}{\second} (superando los \qty{1.48}{\second} simulados) y sostuvo la vertical invertida durante \qty{14.021}{\second} con cero caídas, alcanzando $\mathrm{RMS}(\tilde{\alpha})=\qty{1.99}{\degree}$ y $u_{\mathrm{RMS}}=\qty{0.73}{\volt}$ en el régimen completo (\qty{1.46}{\degree} y \qty{0.57}{\volt} en la ventana óptima de \qty{5}{\second}; hasta \qty{1.35}{\degree} y \qty{0.50}{\volt} en ensayos extendidos), con rechazo validado ante perturbaciones impulsionales (retorno a la banda de $\pm\qty{3}{\degree}$ en menos de \qty{0.9}{\second} en simulación) y rearme autónomo mediante \texttt{RECOVER\_HOME}.
\end{enumerate}

\subsection{Recomendaciones Prácticas y Líneas de Trabajo Futuro}\label{subsec:recomendaciones_futuro}

Para la replicación y el diseño embebido de sistemas mecatrónicos subactuados sobre hardware de bajo costo, se establecen cuatro recomendaciones prácticas derivadas de la evidencia experimental: \textit{(i)}~ubicar el sensor acimutal en el eje del rotor previo a la reductora planetaria cuando se utilicen cajas de alta relación, multiplicando mecánicamente la resolución efectiva sin recurrir a sensores ópticos; \textit{(ii)}~dedicar buses Soft-$\text{I}^2\text{C}$ independientes a sensores con dirección esclava fija para eliminar latencias de multiplexación y bloqueos de periféricos; \textit{(iii)}~implementar los filtros \textit{biquad} selectivos en precisión doble de 64~bits con estructura TDF-II para preservar la ubicación exacta de las muescas cuando los polos se aproximan al círculo unitario, empleando siempre compensación continua por tramos de la zona muerta y evitando integradores competitivos en el eje no actuado; y \textit{(iv)}~aislar el lazo determinista de \qty{1.0}{\kilo\hertz} dentro de la interrupción de temporizador de alta prioridad, desacoplando por completo la transmisión CSV por \texttt{DMA1} a \qty{100}{\hertz} y el renderizado del gemelo digital~3D a \qty{30}{\hertz}.

Con fundamento en las limitaciones físicas y fuentes de no idealidad analizadas en la Sección~\ref{sec:discusiones} ---caída de saturación y asimetría Darlington del puente~H, cuantización angular distal de 12~bits, holgura del tren epicicloidal, fricción de Coulomb en el rodamiento y par elástico del cableado móvil---, se proponen cuatro líneas de investigación futura:

\begin{itemize}
    \item \textbf{Etapa de potencia MOSFET sincrónica:} Reemplazar el puente~H bipolar L298N por un driver MOSFET sincrónico con baja resistencia de encendido inferior a \qty{10}{\milli\ohm}, suprimiendo la caída de conducción cercana a \qty{2.0}{\volt} y la asimetría de zona muerta para extender la linealidad de actuación a baja tensión y reducir las pérdidas térmicas.
    \item \textbf{Encoders ópticos en cuadratura por hardware (\texttt{TIM2}/\texttt{TIM5}):} Sustituir la lectura magnética distal por encoders ópticos incrementales acoplados directamente a los contadores en cuadratura \texttt{TIM2} y \texttt{TIM5} del microcontrolador STM32F446RE, elevando la resolución angular pendular sin carga adicional de procesador y atenuando los picos de ruido derivativo a \qty{1.0}{\kilo\hertz}.
    \item \textbf{Control Predictivo No Lineal (NMPC) en la FPU Cortex-M4F:} Aprovechar la unidad de punto flotante por hardware del procesador para ejecutar en tiempo real un esquema NMPC que unifique el bombeo energético de levantamiento y la estabilización invertida en un único optimizador con restricciones explícitas de tensión, pendiente y carrera acimutal.
    \item \textbf{Observador de estados y perturbaciones por Filtro de Kalman Extendido (EKF):} Integrar un observador no lineal EKF para estimar velocidades angulares limpias con mínimo retardo de fase y reconstruir en línea los pares parásitos inducidos por la fricción seca distal, la holgura de engranajes y la rigidez elástica del cableado con miras a su compensación adaptativa prealimentada.
\end{itemize}

Finalmente, todos los recursos de ingeniería generados en este trabajo ---incluyendo los archivos CAD paramétricos, las placas de manufactura aditiva, el firmware determinista de la STM32F446RE, los registros experimentales de banco, el gemelo digital~3D en MATLAB, el simulador cinemático interactivo~3D (\texttt{simulacion cinematica/\allowbreak index.html}) y las fuentes del artículo--- se encuentran disponibles abiertamente en \url{https://github.com/Nocudo/Pendulo-Furuta-Control-Digital} para respaldar la verificación independiente y el desarrollo de estas extensiones futuras.
